\section{Simulated Gradient} \label{sec:modification}
\simulator{} adapts backpropagation to compute gradients 
~(\Cref{fig:general_structure_simulator}).
We aim to train a capable (i.e., pre-trained) auxiliary model for just a few steps, so high precision gradients may be unnecessary. Instead, \simulator{} performs an approximate backpropagation.  
\simulator{} then uses this \emph{simulated gradient} to update the auxiliary model.
Prior works computed similar approximate gradients in hopes of more faithfully modeling neurobiology~\citep{scellier2017equilibrium,hinton2022forwardforward} or improving the efficiency of training models~\citep{hu2021lora,malladi2023finetuning}.
\iffalse In this section, we describe how \simulator{} performs an approximate backpropagation to compute a \emph{simulated gradient} of the auxiliary model. Due to the complexity of transformer modules, computing {\em exact} gradients of a transformer auxiliary model can require a huge simulator model. Our experiments show that the simulated gradient is a reasonable approximation. \fi 
\iffalse that makes \simulator{} much more expressive than prior constructions despite being orders of magnitude smaller and defer the complete formal construction to the appendix.\fi 
We note that the approximations in the simulated gradients can be made stronger at the cost of enlarging \simulator{}.
Indeed, one could construct a simulator to \emph{exactly} perform the procedure outlined in \textsection{\ref{sec:design}}, though it would be orders of magnitude larger than \simulator{}.
For brevity's sake, we focus on the key approximations and design choices and defer formal details to the appendix.

\subsection{First-order approximations} \label{sec:first_order_approx}
We use first-order approximations of gradients to backpropagate through 
%training dynamics to update the parameters of 
the layer normalization layer.\footnote{We discuss a layer normalization layer $\layernorm$ without scale and bias parameters, but \cref{sec:LNappendix} contains a general construction.} 
It normalizes the input using its mean and standard deviation across the input dimensions. Since the operation is token-wise, we can consider a single position $t$ without loss of generality.

\begin{definition}[Layer normalization]
     A layer normalization layer $\layernorm$  takes input $\vx\in\RR^{\aux{D}}$ and outputs $\vy = (\vx - \mu) / \sigma$, where $\mu$ and $\sigma$ denote its mean and standard deviation.
\end{definition}
\textbf{High precision gradients:} 
Formally, for input-output pair $(\vx, \vy)$, we can compute the  gradients $\partial_{\vy}$, $\partial_{\vx}$ with  chain rule:
\begin{gather}
    \partial_{\vx} = \left(\frac{\partial \layernorm (\vx)}{\partial \vx} \right)^{\top} \losspartial{\vy}  \nonumber \\= \frac{1}{\sigma} \left(  \langle \losspartial{\vy}, \vy \rangle \vy + \losspartial{\vy} - \frac{1}{\aux{D}} \sum_{i=1}^{\aux{D}} \losspartial{y_i} \right). \label{eq:ln_informal}
\end{gather}
\textbf{Inefficiency of exact computation:} 
A \simulator{} layer simulating backpropagation through an auxiliary's layer normalization layer receives  $\losspartial{{\vy}_t}$ and $\vx_t$ in its input embeddings. We go through the exact gradient and why it is inefficient.

For exact computation one could first compute $\vy_t$ using a normalization layer and store in the embeddings. However, inefficiency arises from computing the term  $\langle \losspartial{{{\vy}_t}}, \vy_t \rangle \vy_t$.
To calculate $\langle \losspartial{{{\vy}_t}}, \vy_t \rangle \vy_t$ at each token position $t$, we could either: (1) use a two-layer MLP that focuses on each token separately, or (2) a single self-attention module to treat the operation as a sequence-to-sequence task. 

For (1) we could initially compute $\langle \losspartial{{\vy}_t}, \vy_t \rangle$ via an MLP, followed by computation of $\langle \losspartial{{\vy}_t}, \vy_t \rangle \vy_t$ using another MLP. The element-wise multiplication in embeddings would be facilitated with a nonlinear activation function like \gelu{}~\citep{akyurek2022learning} (refer to thm. \ref{thm:gelu_multiplication} for details). However, this approach would need  substantial number of simulator parameters to represent the MLPs.

\looseness-1Alternatively, we could  use a single self-attention module. Constructing such a module would require careful engineering to make sure the input tokens only attend to themselves while keeping an attention score of $0$ to others. If we used a linear attention, we would need to space out the gradient $\losspartial{{\vy}_t}$ and $\vx_t$ in each position $t$, such that the attention score is $0$ between different tokens. This would require an embedding dimension proportional to the context length. 
On the other hand, if we used a softmax attention module, we would need an additional superfluous token in the sequence. Then, a token at position $t$ would attend to itself with attention $\langle \partial{\vy}_t, \vy_t\rangle$ and to the extra token with an attention score of $1-\langle \partial{\vy}_t, \vy_t\rangle$. The extra token would return a value vector $0$. \sm{Rephrase the prior sentence...I can't understand how this would work} To avoid such inefficiency, we opt for a first-order approximation instead.


\textbf{Efficient approximation:} 
Instead of explicitly computing each term in the chain rule of $\left(\frac{\partial \layernorm (\vx)}{\partial \vx} \right)^{\top} \losspartial{\vy}$ in Eq. \ref{eq:ln_informal}, we instead use a first order Taylor expansion of $\layernorm$.
\begin{align*}
    \layernorm(\vx + \epsilon \losspartial{\vy}) = \layernorm(\vx) + \epsilon \left(\frac{\partial \layernorm (\vx)}{\partial \vx} \right) \losspartial{\vy} + \mathcal{O}(\epsilon^2).
\end{align*}
Rearranging allows us to write
\begin{align*}
    \left(\frac{\partial \layernorm (\vx)}{\partial \vx} \right) \losspartial{\vy} =  \frac{1}{\epsilon} \left( \layernorm (\vx + \epsilon \partial_{\vy}) - \layernorm (\vx) \right) + \mathcal{O}(\epsilon).
\end{align*}
Similar to the computation of Eq. \ref{eq:ln_informal}, we can show  
\begin{align*}
    \frac{\partial \layernorm (\vx)}{\partial \vx} = \frac{1}{\sigma} \left( (1-{\aux{D}}^{-1})\mI - \layernorm(\vx) \layernorm(\vx)^{\top} \right).
\end{align*}
Because  $\partial \layernorm (\vx) / \partial \vx$ is symmetric\footnote{For a linear function $f$ with matrix $\mW$, $\frac{\partial f (\vx)}{\partial \vx} = \mW$. Since $\mW$ may not be a symmetric matrix, this method can't be generally applied to approximately backpropagate linear layers or causal self-attention layers.}, we can write
\begin{align*}
    \left(\frac{\partial \layernorm (\vx)}{\partial \vx} \right)^{\top} \losspartial{\vy} & =  \left(\frac{\partial \layernorm (\vx)}{\partial \vx} \right) \losspartial{\vy} \\& =  \frac{1}{\epsilon} \left( \layernorm (\vx + \epsilon \partial_{\vy}) - \layernorm (\vx) \right) + \mathcal{O}(\epsilon).
\end{align*}



Then, ignoring the small error term, we can use just two linear layers, separated by a normalization layer, to simulate the approximation. 


\subsection{Fuzzy backpropagation via stop gradients} \label{sec:stop_gradients}
Self-attention is inherently quadratic, because it uses the keys and queries to compute attention scores between every possible pair of tokens in the sequence.
These scores then linearly combine the value vectors (see def. \ref{def:self-attn_auxiliary_single}). 
Computing the gradient exactly is thus a very complex operation.
Instead, we stop the gradient computation through attention scores in the self-attention layer. 
For similar reasons, we only update the value parameter in the self-attention module.

\textbf{Gradient backpropagation:}
For an input, output sequence pair $\{\vy_t\}, \{\vy_t\}$, if $\{\vq_t, \vk_t, \vv_t\}$ denote the intermediate query, key, value vectors, on gradients $\{ \partial_{\vy_t} \}$, $\{ \partial_{\vx_t} \}$ is given via the chain rule:
\begin{align}
    \losspartial{\vx_t} &= \mQ^{\top} \losspartial{\vq_t} + \mK^{\top} \losspartial{\vk_t} + \mV^{\top} \losspartial{\vv_t}.
    \label{eq:attn_backprop_informal}
\end{align}
Here, $\mV, \mK, \mQ$ denote the query, key, and value matrices. 

\textbf{Inefficiency in exact computation:} 
Here, we demonstrate that simulating computation of the three terms in Eq. \ref{eq:attn_backprop_informal} is inefficient, because
$\losspartial{\vq_t}, \losspartial{\vk_t}$ depend on the derivatives w.r.t. the attention scores. 
As an example, we focus on $\losspartial{\vk_t}$:
\begin{align*}
    \losspartial{\vk_t} = \sum_j a_{t, j} ( (\losspartial{\vy_t})^{\top} \vv_j ) (\vk_j - \sum_{j'} a_{t, j'} \vk_{j'}).
\end{align*}

Computing this term would require us at least 2 self-attention layers and an MLP layer. The first attention layer would compute $(\losspartial{\vy_t})^{\top} \vv_j$ for different token pairs, similar to the forward simulation of a linear layer with linear attention (\textsection{\ref{sec:exposition_linear_forward}}). These would be then multiplied to the pair-wise attention scores $a_{t, j}$  with an MLP to compute $a_{t, j} ( (\losspartial{\vy_t})^{\top} \vv_j )$, with elementwise product would be facilitated by GeLU non-linearity (thm. \ref{thm:gelu_multiplication}). 
These would be finally used by an attention layer to combine the different key vectors. A similar simulation would be necessary to compute $\losspartial{\vq_t}$.




\textbf{Stop gradients through query and key vectors:}
In order to reduce the necessary resources, we ignore the query and key gradients in Eq. \ref{eq:attn_backprop_informal}. When we ignore these gradient components,  $\{ \partial_{\vx_t} \}$ can be simplified as
\begin{align}
    \losspartial{\vx_t} \approx \mV^{\top} \losspartial{\vv_t} = \mV^{\top} \sum_{j} a_{j, t} \losspartial{\vy_t}.
    \label{eq:approx_attn_gradient}
\end{align}
A single self-attention layer can compute this by using the attention scores to combine the token-wise gradients.

\textbf{Why won't it hurt performance?} Estimating $\partial_{\vx_t}$ as described is motivated by recent work \citep{malladi2023finetuning} showing that fuzzy gradient estimates don't adversely affect fine-tuning of pre-trained models. Furthermore, we theoretically show that when the attention head for each position pays a lot of attention to a single token (i.e., behaves like hard attention \citep{perez2021attention}), the approximate gradient in Eq. \ref{eq:approx_attn_gradient} is entry-wise close to the true gradients (thm. \ref{thm:attn_backprop}).

The other approximation is to update only the value parameters $\mV$ of the auxiliary model (\textsection{\ref{sec:self-attnt-backprop_appendix}}). 
This is motivated by parameter efficient fine-tuning methods like \lora{} \cite{hu2021lora} and \ia{} \cite{liu2022few}, which restrict the expressivity of the gradient updates without degrading the quality of the resulting model.
We similarly show in the next section that the simulated gradients in \simulator{} can effectively tune large pre-trained transformers.
